\SourceLeaf{219}%
\noindent nach der Annahme, dass aber der Satz

\[ c\;[\text{1ste Figur}] = c\;[\text{2te Figur}] \]

\noindent dann bewiesen werden müsse?

Mir scheint bei der Willkürlichkeit der Winkel dieser Beweis nicht nothwendig.

Diess sind die Grundzüge des Beweises und ich erwarte Ihre Entscheidung. Ich füge nur, um meinen Beweis zu rechtfertigen, hinzu, dass freilich durch die zweite Operation das Dreieck $ABC$ verschwindet, aber \textls{nicht} die Winkel des Dreiecks. Wie die Linien auch liegen, so ist immer

\[ J\hat{B}H = \beta, \qquad G\hat{C}F = \gamma, \qquad D\hat{A}E = \alpha \]

\noindent im endlichen, so wie im verschwindenden Dreiecke, mithin die Summe

\[ J\hat{A}H + G\hat{A}F + D\hat{A}E \]

\noindent immer gleich der Summe der Winkel eines geradlinigen Dreiecks.

Soll man also den Satz von einem beliebigen Dreiecke (dessen Winkel $A$, $B$, $C$) beweisen, so zieht man die Linien $DG$, $EH$, so dass

\[ \alpha = A, \]

\noindent man nimmt ferner den Winkel $J\hat{A}H = B$, $G\hat{A}F = C$.

Ist dann $JAF$ keine gerade, sondern eine gebrochene Linie $JAF'$, so ist freilich der Winkel $c$ dadurch um $dc$ kleiner, der Winkel $b$ aber um eben so viel grösser geworden, mithin ihre Summe unverändert geblieben, oder wir haben, was zur Stringenz des Beweises gehört:

\[ b + c\;[\text{Fig. 1}] = b + c\;[\text{Fig. 2}].\} \]

\bigskip
\hrule
\bigskip

\Needspace{7\baselineskip}
\begin{center}
\textsc{Gauss} an \textsc{Schumacher}. Göttingen, 17. Mai 1831.
\end{center}

\mbox{.\ .\ .\ .\ .\ .\ .} Bei dem, was Sie über die Parallellinien schreiben, haben Sie, genau besehen, in Ihren Syllogismen einen Zwischensatz gebraucht, ohne ihn ausdrücklich auszusprechen, der so lauten müsste:
